\documentclass[11pt]{article}
\usepackage[letterpaper,margin=1in]{geometry}
\usepackage{amsmath,amssymb,amsthm}
\usepackage[colorlinks=true,urlcolor=blue,linkcolor=blue,citecolor=blue]{hyperref}
\hypersetup{pdftitle={Antichains in locally finite trees without countable choice},pdfauthor={Wentao Li}}
\newtheorem{theorem}{Theorem}
\newtheorem{lemma}[theorem]{Lemma}
\newcommand{\II}{\mathcal I}
\newcommand{\UU}{\mathcal U}
\newcommand{\FF}{\mathbb F_2}
\newcommand{\supp}{\operatorname{supp}}
\title{Antichains in locally finite trees without countable choice}
\author{Wentao Li}
\date{Research draft prepared with assistance from OpenAI GPT-6.2 Astra}
\begin{document}
\maketitle
\begin{abstract}
Relative to the consistency of ZF, every uncountable locally finite rooted tree can contain an uncountable antichain while a countable sequence of pairs has no choice function. This gives simultaneous negative answers to the two implications in Question~6.6 of Banerjee and Gopaulsingh, in pure ZF as well as with atoms. A selective ultrafilter governs the allowed supports of a permutation model. Linear characters of the coordinate flips detect exactly which supports can enumerate a vertex orbit. Separately, in Pincus's Model X even well-orderable families of nonempty finite sets have choice functions, answering their Question~6.7.
\end{abstract}

\section{The questions}
Countable means admitting an injection into $\omega$. A rooted tree is a connected acyclic graph with a distinguished vertex. Its order is ancestry. In a locally finite rooted tree, each vertex has finitely many children. All its levels are finite, by ordinary finite induction in ZF. An $\omega$-tree additionally has points at every finite height (up to a shift of one in the numbering). Banerjee and Gopaulsingh \cite[Definitions 2.2, 2.4 and Question 6.6]{bg} consider two assertions. An $\omega$-tree is countable if all its chains are finite and its antichains countable, or, respectively, if all its chains are countable and its antichains finite. They ask whether either assertion implies countable choice for finite sets. They prove the implications in the opposite direction.
\begin{theorem}\label{main}
Relative to the consistency of ZF there is a model of pure ZF satisfying both
\begin{enumerate}
\item Every uncountable locally finite rooted tree has an uncountable antichain.
\item Some sequence of two-element sets has no choice function.
\end{enumerate}
There is also a permutation model of ZFA with these properties.
\end{theorem}
The first clause implies both assertions in the question. In the model there really is an uncountable $\omega$-tree. For example the tree of finite initial choices from the sequence of pairs cannot be countable, since a well-order of its vertices would canonically choose an extension at each step and thereby give a branch.

\section{A model with restricted sets of coordinates}
Begin in a ground universe with choice and CH and a countable sequence of disjoint pairs of atoms indexed by $\omega$. This ground theory has the consistency strength of ZF, by constructibility and the usual adjoining of atoms. Fix a nonprincipal Ramsey ultrafilter $\UU$ on $\omega$ in the ground universe. Recall a construction: list, in $\omega_1$ stages, all subsets of $\omega$ and all finite colorings of pairs. Build infinite sets decreasing modulo finite. At each countable limit take an infinite diagonal pseudointersection of earlier sets (finite intersections of the earlier infinite sets are infinite). At a successor restrict to an infinite side of the next subset and an infinite homogeneous subset for the next coloring, by ordinary Ramsey. Their upward closure modulo finite is the desired ultrafilter. Every countable family of its members has a pseudointersection belonging to it, by taking a stage past their witnessing stages.

Let $\II=\{A:\omega\setminus A\in\UU\}$. Given countably many members of $\II$, there is an $F_0\in\II$ containing each of them modulo finite. Let $G=\bigoplus_\omega\FF$ act by flipping the indicated pairs. Only finitely many pairs move in each group element. For $E\in\II$, let $G_E$ consist of the flips that avoid $E$. Use the normal filter of subgroups generated by the $G_E$ to form the hereditary permutation model. Each pair and the whole indexed sequence are internal. A selector supported by $E$ would have to be fixed by the single flip at any $j\notin E$, an impossibility. In what follows, constructions and choices made in the ground universe are brought into the model by giving a support and checking that their members are internal.

\section{Finite orbits in a tree}
Fix an internal locally finite rooted tree and a support $E$ for its entire structure. Let $K=G_E$. The tree is externally countable as a union of finite levels. Each $K$-orbit $O$ of vertices lies in one finite level. Its stabilizer $L_O$ is the same for all its points, since $K$ is abelian. The finite quotient $K/L_O$ acts regularly on $O$. Define
\[
 D(O)=\{d\in K^*:d|_{L_O}=0\},\qquad
 \sigma(O)=\bigcup_{d\in D(O)}\supp(d).
\]
Here $K^*$ is the ground algebraic dual. Each character is identified with its coefficient function on $\omega\setminus E$ because flips have finite support. $D(O)$ is the dual of the finite-dimensional $\FF$-space $K/L_O$. It separates the cosets of $L_O$. For every $M\subseteq\omega$ this proves the criterion
\begin{equation}\label{criterion}
 G_{E\cup M}\text{ fixes one (and every) point of }O
 \quad\Longleftrightarrow\quad \sigma(O)\subseteq M .
\end{equation}
Indeed, the individual coordinate flips outside $E\cup M$ generate that subgroup. A support of any vertex of $O$ consequently contains $\sigma(O)$, so $\sigma(O)\in\II$.

The orbit of a parent is the unique parent of the child's orbit in the quotient tree. Along ancestry the subgroups $L_O$ decrease, and $D(O)$ and $\sigma(O)$ increase. Put $\Sigma=\bigcup_O\sigma(O)$. If $\Sigma\in\II$, fixing $E\cup\Sigma$ fixes every vertex, so a ground injection of the tree into $\omega$ belongs to the model. Conversely, the support $M$ of an internal injection into $\omega$ must fix each vertex under $G_{E\cup M}$, giving $\Sigma\subseteq M$ by \eqref{criterion}. Thus the tree is internally uncountable precisely when $\Sigma\in\UU$.

Assume this last condition. For each $j\in\Sigma$ choose externally an orbit $O(j)$ with $j\in\sigma(O(j))$. Color pairs from $\Sigma$ according as their two orbits are equal, distinct comparable, or incomparable in the quotient tree. Extend arbitrarily to $\omega$ and apply the Ramsey property. Some $J\in\UU$, $J\subseteq\Sigma$, is homogeneous. Equality would put $J$ inside one small $\sigma(O)$. For incomparable orbits the union of their vertices is a $K$-invariant antichain, so it is internal with support $E$. An internal injection of it into $\omega$ with support $M$ would give $J\subseteq M\in\II$ by \eqref{criterion}. This proves the desired conclusion in that case.

\section{When the quotient orbits form a chain}
In the remaining case the chosen orbits are distinct and comparable. They have unbounded heights since each level is finite. Their ancestors give a sequence $O_n$ at consecutive levels starting with the root. Each parent map $O_n\to O_{n-1}$ is onto, by equivariance and transitivity. Choose a ground ray $t_n\in O_n$ along these maps. This ray need not be internal. Write $D_n=D(O_n)$ and $D=\bigcup_n D_n$. Its coefficient support contains $J$. Hence $D$ is infinite-dimensional: a finite basis would have small total support. It is countable in the ground universe.

\begin{lemma}\label{characters}
For some $F\in\II$ the subspace of $D$ consisting of characters supported on $F$ is infinite-dimensional.
\end{lemma}
\begin{proof}
Choose $F_0\in\II$ almost containing the support of every $d\in D$, using the pseudointersection property above. Restriction to $\omega\setminus F_0$ maps $D$ into the space of finitely supported functions there. An infinite-dimensional kernel gives the result with $F_0$. Otherwise the image has infinite dimension. Choose nonzero image vectors with pairwise disjoint finite supports: the subspace vanishing on the union of finitely many previous supports has finite codimension in the image. Take preimages in $D$. At least one of the two disjoint unions of the even and odd supports is outside $\UU$. Adjoin that union to $F_0$ to obtain $F\in\II$. The corresponding preimages are supported on $F$ and are independent since their images are independent.
\end{proof}

Put $H=G_{E\cup F}$. The number of $H$-orbits in $O_n$ equals
\[
 [K:H+L_{O_n}]=2^{\dim\{d\in D_n:d|_H=0\}}.
\]
This follows by taking the dual of the finite quotient. Lemma~\ref{characters} makes these nondecreasing numbers unbounded. For each $n$ where the number strictly increases over level $n-1$, take all the vertices in $O_n\setminus Ht_n$ whose parent lies in $Ht_{n-1}$. There are such vertices: between the finite quotients by $H+L_{O_n}$ and $H+L_{O_{n-1}}$ the parent map is a surjective map of coset spaces, all fibers having the same size. At a strict increase each fiber has at least two points.

The union of these sets is $H$-invariant and is internal. It is an antichain. A vertex taken at an earlier jump is outside $Ht_n$ at its own level $n$, whereas any vertex taken later has its ancestor at level $n$ in $Ht_n$, by equivariance along the ray. Different vertices at one level are incomparable. A support $M\in\II$ for an injection of this antichain into $\omega$ would contain $\sigma(O_n)$ at every jump by \eqref{criterion}, since it fixes individually the selected vertices under $G_{E\cup M}$. The jump levels are unbounded, so their increasing supports cover $J$. This is impossible. The proof in ZFA is complete.

\section{Transfer to pure sets}
We use Pincus's transfer theorem for a finite conjunction of injectively boundable statements \cite[Theorem~3A3]{pincus}; see also \cite[Theorem~2.6]{bg}. We give the bounds, since the universal assertion quantifies over all tree domains. Write $|X|^-$ for the supremum of the well-ordered cardinals that inject into $X$. Every well-orderable subset $W$ of a locally finite rooted tree injects into $\omega\times\omega$: fix one well-order of $W$ and send a point to its level and its index, in that order, within the finite intersection of $W$ with the level. Hence any vertex domain of such a tree satisfies $|X|^-\le\omega$ in ZFA.

Let $\Psi(X)$ say that every locally finite rooted tree on $X$ either admits an injection of $X$ into $\omega$ or has an antichain admitting no injection into $\omega$. This is a boundable formula. Indeed graph relations and roots, finite paths, subsets, the level witnesses in $\omega$, and maps from subsets of $X$ into $\omega$ are all coded inside $P^{\omega+20}(X)$ under the usual increasing power-set iteration. Local finiteness can be tested by a bijection of each successor set with a finite ordinal. All quantifiers may be restricted to these codes, so truth is absolute to the indicated stage. The tree assertion is therefore equivalent to
\[
  \forall X\ (|X|^-\le\omega\ \longrightarrow\ \Psi(X)).
\]
This is the defining injectively boundable form with no parameters and the boundable constant $\omega$ \cite[Section~5]{mad}. Failure of countable choice for pairs is boundable: existentially specify a set of points $A$, then use a sequence of pairs from $A$ and quantify over all selectors on $\omega$ inside $P^{\omega+20}(A)$. Thus the conjunction transfers. This finishes Theorem~\ref{main}.

\section{Choice in Pincus's Model X}
Question~6.7 of \cite{bg} concerns a different, specified permutation model. In Model X (model $\mathcal N_{34}$ in Howard--Rubin \cite[p.~207]{hr}) there is a pair of atoms over each rational. The group acts by increasing automorphisms of the rational order, lifted without swaps, together with arbitrary \emph{finite} products of individual pair swaps. Supports are finite sets of atoms, and they can be enlarged to finite unions of pairs.
\begin{theorem}\label{modelx}
In Pincus's Model X, the union of a well-orderable family of well-orderable sets is well-orderable. In particular a well-orderable family of nonempty finite sets admits a choice function.
\end{theorem}
\begin{proof}
For a finite set $E$ of rational positions, let $K$ fix both atoms over each position of $E$. We first prove that any normal subgroup $N$ of $K$ containing the pointwise stabilizer of finitely many further pairs is all of $K$. Denote their positions by $D$. The kernel assumption in this statement is pointwise fixation on both $E$ and $D$.

Every single swap at a position outside $E$ belongs to $N$: in its component interval of $\mathbb Q\setminus E$ move the position away from $D$ by an order automorphism fixing $E$, then use normality. Consider next an increasing automorphism supported inside one bounded closed subinterval of one component, lifted without swaps. Some conjugate has support disjoint from $D$, so it too lies in $N$. Here the required moves exist by homogeneity of the countable dense linear order, or by extending an increasing map of finitely many rationals piecewise linearly. Finally, any increasing automorphism fixing $E$ agrees on $D$ with a product of such bounded-interval automorphisms. In each of the finitely many components containing points of $D\setminus E$ enclose those points and their images in a closed subinterval with rational endpoints in the component. Interpolate increasingly between the prescribed images and fix those endpoints and everything outside. The remaining factor fixes $E\cup D$ pointwise. This proves the subgroup assertion for increasing lifts and, using the swaps, for all of $K$.

Fix a support $E$ for a well-ordered indexed family $(X_i)$ and its index order. $K$ fixes each index and preserves each $X_i$. For each $i$ the internal existence of a well-order of $X_i$ gives a finite additional support for one such well-order. Fixation of that support together with $E$ fixes the well-order, hence fixes all the points of $X_i$ individually. The kernel of the action of $K$ on $X_i$ is a normal subgroup satisfying the assertion just proved. Thus $K$ fixes every point of every $X_i$. One ground well-order of their union is consequently hereditarily supported by $E$ and belongs to the model. Taking least elements proves the choice assertion.
\end{proof}
The same argument applies to any family that admits a well-order in the model, by fixing one such order first. For a single sequence of pairs, the difference between these two models is the action on the indices. In the first model each index is fixed and the independent flips can survive in local quotients. In Model X the order automorphisms move rational positions away from every finite support.

\paragraph{Relation to earlier work.} Hall and Shelah \cite{hs} construct permutation models separating partial from full choice for a countable family of sets of prime size. Banerjee and Gopaulsingh \cite[Remark~4.5]{bg} give a tree violating both tree-countability principles in the second Fraenkel model. Theorem~\ref{main} instead verifies both universal tree assertions while a sequence of pairs has no selector. Tachtsis \cite{t24} and Howard--Tachtsis \cite{ht26} resolve other questions about posets and several versions of the graph partition theorem in the same 2023 paper, including its Question~6.5. Model X and its tree chain/antichain property are due to Pincus. Theorem~\ref{modelx} concerns its well-orderable unions. Both arguments and their priority await independent review.
\begin{thebibliography}{10}
\bibitem{bg} A. Banerjee and A. Gopaulsingh. \href{https://doi.org/10.4064/ba221221-6-6}{On the Erdős--Dushnik--Miller theorem without AC}. Bull. Polish Acad. Sci. Math. 71 (2023), 1--21. Questions 6.6 and 6.7.
\bibitem{pincus} D. Pincus. \emph{Zermelo--Fraenkel consistency results by Fraenkel--Mostowski methods}. J. Symbolic Logic 37 (1972), 721--743.
\bibitem{mad} E. Tachtsis. \emph{The Boolean prime ideal theorem does not imply the extension of almost disjoint families to MAD families}. Bull. Polish Acad. Sci. Math. 68 (2020), Section~5.
\bibitem{hr} P. Howard and J. Rubin. \emph{Consequences of the Axiom of Choice}. AMS, Mathematical Surveys and Monographs 59 (1998), Model 34 and Note 105.
\bibitem{hs} E. Hall and S. Shelah. \href{https://arxiv.org/abs/0808.0535}{Partial choice functions for families of finite sets}. Fund. Math. 220 (2013).
\bibitem{t24} E. Tachtsis. \href{https://doi.org/10.1007/s00605-023-01933-z}{On the deductive strength of the Erdős--Dushnik--Miller theorem and two order-theoretic principles}. Monatsh. Math. 203 (2024), 677--693.
\bibitem{ht26} P. Howard and E. Tachtsis. \href{https://doi.org/10.4064/fm250514-20-7}{Three forms of the Erdős--Dushnik--Miller theorem}. Fund. Math. 272 (2026), 171--203.
\end{thebibliography}
\end{document}
